\subsection{Analyse}

Voici un tableau récapitulatif des résultats issus des différentes 
manipulations réalisées. 

\begin{table}[h!]
    \centering
    \begin{tabular}{|l|c|c|c|c|c|c|}
        \hline
        & Image $\infty$ & Autocollim. & Silbermann & Bessel & Pts Conj. & Badal \\
        \hline
        $f'$ (cm) & 9,5 & 10,2 & 10,0 & 9,99 & 10,1 & -10,0 \\
        \hline
        $\Delta f'$ (cm) & 0,2 & 0,4 & 0,1 & 0,05 & 0,2 & 0,05 \\
        \hline
        $\Delta f' / f'$ (\%) & 2,1 \% & 3,9 \% & 1,0 \% & 0,5 \% & 2,0 \% & 0,5 \% \\
        \hline
    \end{tabular}
    \caption{Synthèse globale des mesures focométriques pour $L_1$ et $L'$.}
\end{table}

\subsection{Discussion}

L'analyse de la synthèse des résultats met en évidence une corrélation directe 
entre la méthodologie employée et la précision finale obtenue.
Les méthodes de \textbf{Bessel} et de \textbf{Badal} présentent 
les incertitudes relatives les plus faibles ($0,5\%$).
Ce résultat remarquable s'explique par le fait que ces protocoles 
reposent quasi exclusivement sur la mesure de déplacements mécaniques 
effectués sur le banc optique.
Cette approche permet de minimiser l'impact de l'appréciation visuelle 
au profit d'une lecture instrumentale plus fiable.

En revanche,
les méthodes telles que l'autocollimation ou les points conjugués 
sont fortement limitées par la détermination subjective 
de la netteté de l'image sur l'écran.
Plusieurs facteurs environnementaux et humains 
viennent dégrader la précision de ces mesures :
la luminosité ambiante de la salle qui réduit le contraste,
la fatigue visuelle de l'expérimentateur lors des essais répétés,
ainsi que la qualité intrinsèque du faisceau qui peut paraître « baveux » 
en raison des aberrations optiques.

Il est important de souligner que les incertitudes présentées 
dans notre tableau pourraient être sous-évaluées.
Une approche plus rigoureuse consisterait à intégrer systématiquement 
une \textit{incertitude de latitude de mise au point} 
liée à la profondeur de champ perçue par l'œil.
L'inclusion de ce paramètre augmenterait mécaniquement les incertitudes absolues,
particulièrement pour les méthodes dépendant d'une observation directe.
Cependant,
cette analyse plus fine ne ferait que renforcer notre conclusion principale :
la supériorité des méthodes différentielles et de déplacement 
par rapport aux méthodes de visée directe.